\section{Evaluation}
\label{sec:evaluation}
We ask four questions. \textbf{RQ1} tests whether legal repeated-token inputs produce the predicted small orbit and identify an equivalent demixer. \textbf{RQ2} measures token-multiset recovery across models and configurations. \textbf{RQ3} tests the official entry point and the effect of secret lifetimes. \textbf{RQ4} asks whether the recovered content supports a concrete, bounded privacy task and whether shortening the secret lifetime interrupts transfer.

\subsection{Experimental methodology}
\label{sec:methodology}
The main suite uses the official \scheme{} implementation at commit \texttt{6b40f36}. We do not modify its cache transformation. Model forward passes use FP32 weights on CPU, and the cloak stores/transforms cache blocks in bf16, matching the implementation default. Experiments ran on an AMD EPYC 9554 host. Table~\ref{tab:breadth} reports ten distinct public models from the Qwen, Llama, OLMo, Phi, and TinyLlama families~\cite{qwen25,llama3,olmo,phi3,tinyllama} and twelve configurations; Qwen2.5-0.5B additionally covers block sizes 8 and 32. Checkpoint revisions are listed in Appendix~\ref{app:repro}.

For every configuration, the attacker submits text through the unmodified tokenizer. Any beginning-of-sequence or other prefix token remains in the model input. Calibration uses only complete blocks inside the final constant-token run; content recovery is scored on complete blocks of the target texts. Calibration requires all $b$ states to be observed for a head. Every head in the main suite calibrated, so Table~\ref{tab:breadth} includes all heads.

Five target texts cover English prose, Python, Chinese, dialogue, and mathematical prose. Each is evaluated under two fresh block-permutation seeds. The primary metric is the multiset hit count in~\eqref{eq:metric}, summed over complete blocks and calibrated heads. Thus a ``head--token opportunity'' is one token position scored through one KV head. Opportunities within a document, across heads, and across seeds are correlated. We report their raw descriptive total and per-configuration rates; we do not use a token-level binomial interval as evidence of generalization to independent requests.

The main-suite fresh-secret control applies a demixer calibrated under one complete configuration to a cache produced under another, scoring head 0 on the English prose target. Its denominator therefore differs from the multi-head, five-text primary metric. Ground-truth secrets are used only for diagnostics such as parameter similarity. The attack routine receives the public model, its own protected calibration cache, and protected target blocks.

\paragraph{Configuration and input accounting}
The main runner estimates the defender's magnitude thresholds from eight fixed natural-language sentences stored separately from the five target texts. It generates one secret configuration per model/block-size cell with seed 42, then evaluates two target-cache permutation seeds, 5001 and 5002. These two seeds change the block permutations; they do not constitute two independently generated secret configurations. The additional multi-epoch study below varies the complete configuration and answers a different question.

Probe lengths are recorded after tokenization. Character repetition can produce different token counts across tokenizers, and the search may overshoot its target length. The table therefore reports the actual submitted token count rather than a nominal character or token budget. Prefix lengths are one token for the recorded Llama and Phi configurations and two for TinyLlama; the remaining main-suite configurations have no prefix before the constant-token run. Prefix-containing blocks and incomplete final probe blocks are excluded from calibration. Target texts are independently truncated to their last complete block, so the reported recovery denominator excludes their trailing partial blocks.

These choices define the unit of each result. The main table aggregates head--token opportunities under the fixed texts and recorded configurations. The probe curve measures repeated calibration trials, the epoch study averages session rates under new configurations, and the privacy study ranks a supplied candidate set. Their denominators are not interchangeable: adding their reported successes into one grand total would mix different questions and levels of aggregation.

\subsection{RQ1: orbit and parameter recovery}
Every evaluated head forms exactly $b$ clusters after prefix blocks are removed: 8, 16, or 32 according to the configuration. The smallest saved between/within cluster-separation ratio is 500.2. Evaluation matches recovered columns of $S$ to true columns by cosine similarity and checks that the matching is bijective. The minimum matched cosine similarity is 0.999996; the largest relative error recorded for $w$ is 0.0354. These evaluation-only diagnostics support the predicted rank-one state geometry.

Figure~\ref{fig:probe} compares the observed two-head calibration rate on Qwen2.5-0.5B with the exact one-head occupancy probability for independent uniform states. Each empirical point contains ten trials under one fixed secret configuration. This diagnostic counts a trial as successful when both heads collect all states and the recorded marker and minimum matched-column cosine similarities each reach 0.999. With 80 blocks (1,280 repeated tokens), 9/10 trials meet this criterion; with 128 blocks (2,048 tokens), all ten do. The small, independently seeded trial sets can yield a non-monotone empirical curve. The comparison supports the coupon-collector scale rather than estimating a deployment success curve.

\begin{figure}[t]
\centering
\includegraphics[width=\columnwidth]{probe_budget.pdf}
\caption{Calibration versus probe length for $b=16$. The empirical series requires full recovery of both heads; the dashed series is the exact one-head occupancy probability. Ten trials are used at each empirical budget.}
\label{fig:probe}
\end{figure}

\subsection{RQ2: content recovery}
Across the twelve configurations in Table~\ref{tab:breadth}, the attack obtains 263,029 hits in 263,904 head--token opportunities, or 99.668\%. Per-configuration recovery ranges from 99.318\% on Phi-3-mini to 100\% on six configurations. Fresh-secret control recovery ranges from 0 to 0.625\%: ten configurations have zero hits, while the $b=8$ and $b=32$ configurations each have one hit. These controls support the role of the recovered secrets in the observed recovery.

\begin{table*}[t]
\centering
\caption{Main-suite recovery from first-layer Value caches. ``Probe'' is the complete tokenizer output length; C/N is the raw multiset hit count. The control uses a fresh secret configuration.}
\label{tab:breadth}
\setlength{\tabcolsep}{3.5pt}
\begin{tabular}{lrrrrrrr}
\toprule
Model & $b$ & Heads & $d$ & Probe & C/N & Rec. (\%) & Control \\
\midrule
Llama-3.2-1B & 16 & 8 & 64 & 3,201 & 20,192/20,224 & 99.84 & 0/176 \\
OLMo-2-1B & 16 & 16 & 128 & 3,200 & 45,043/45,056 & 99.97 & 0/176 \\
Phi-3-mini & 16 & 32 & 96 & 4,080 & 111,872/112,640 & 99.32 & 0/208 \\
Qwen2.5-0.5B & 16 & 2 & 64 & 3,200 & 4,800/4,800 & 100.00 & 0/176 \\
Qwen2.5-0.5B & 32 & 2 & 64 & 20,250 & 4,608/4,608 & 100.00 & 1/160 \\
Qwen2.5-0.5B & 8 & 2 & 64 & 1,280 & 4,896/4,896 & 100.00 & 1/176 \\
Qwen2.5-1.5B & 16 & 2 & 128 & 3,200 & 4,768/4,800 & 99.33 & 0/176 \\
Qwen2.5-3B & 16 & 2 & 128 & 3,200 & 4,800/4,800 & 100.00 & 0/176 \\
Qwen2.5-7B & 16 & 4 & 128 & 3,200 & 9,578/9,600 & 99.77 & 0/176 \\
Qwen3-0.6B & 16 & 8 & 128 & 3,200 & 19,200/19,200 & 100.00 & 0/176 \\
Qwen3-1.7B & 16 & 8 & 128 & 3,200 & 19,200/19,200 & 100.00 & 0/176 \\
TinyLlama-1.1B & 16 & 4 & 64 & 2,033 & 14,072/14,080 & 99.94 & 0/208 \\
\midrule
\multicolumn{5}{r}{Total} & 263,029/263,904 & 99.668 & -- \\
\bottomrule
\end{tabular}
\end{table*}

We separately exercise two implementation variants whose records do not share the main table's denominator. In the fused path, an isolated attacker instance uses the original public model and recovered transform; it does not read the defender's fused weights. Recovery is 100\% on each of the English, Chinese, and Python targets tested. In the optional \texttt{need\_ratio} path ($S$ and $M$ ratio 1.25), the variant-specific recovery succeeds on all five target texts and both permutation seeds. These results establish reachability for the tested variants, not coverage of their full parameter spaces.

\paragraph{Where the misses occur}
The main total contains 875 missed head--token opportunities. Phi-3-mini contributes 768 of them and also has the largest denominator, 112,640, because its first Value layer has 32 KV heads. Its per-configuration recovery remains 99.318\%. This concentration cautions against reading the aggregate as an equal-weight average over models: both the number of heads and the tokenized text lengths affect a configuration's weight.

Table~\ref{tab:textcounts} regroups the same raw counts by target text. It introduces no additional samples. Each row represents one fixed text evaluated repeatedly across configurations, heads, and the two permutation seeds; it is not an estimate over a language or application domain. The breakdown shows that the aggregate is not supplied by one text alone. It cannot establish, for example, that Chinese is easier to recover than English, since vocabulary, length, content, and model weighting all differ.

\begin{table}[t]
\centering
\caption{The main-suite counts regrouped by fixed target text. These are the same correlated evaluation opportunities as Table~\ref{tab:breadth}, not additional experiments.}
\label{tab:textcounts}
\begin{tabular}{lrr}
\toprule
Target text & Matches / opportunities & Misses \\
\midrule
English prose & 33,795 / 33,920 & 125 \\
Python code & 58,587 / 58,784 & 197 \\
Chinese prose & 80,517 / 80,768 & 251 \\
Dialogue & 37,928 / 38,048 & 120 \\
Mathematical prose & 52,202 / 52,384 & 182 \\
\midrule
Total & 263,029 / 263,904 & 875 \\
\bottomrule
\end{tabular}
\end{table}

The retained counters do not identify the cause of each miss. Finite-precision parameter error, marker-hypothesis selection, and dictionary ambiguity are possible explanations, but these aggregates cannot attribute their relative contributions. In particular, a high cosine similarity for recovered parameters does not imply exact token recovery when dictionary margins are small. The conditional margin argument in Section~\ref{sec:analysis} describes that dependency without assigning a measured cause to unretained per-block errors.

The evaluated Value-only attack outputs token multisets. An exploratory ordering script records 9.1\% and 15.9\% position accuracy on 11 blocks with Qwen2.5-0.5B and Qwen2.5-3B, respectively, and no exact blocks. However, it passes unique token keys to beam search and drops multiplicity. Its 6.25\% reference also assumes distinct positions without accounting for repeated token identities. We retain these exploratory observations for transparency but draw no conclusion about the difficulty of order recovery from them.

\subsection{RQ3: entry point and secret lifetime}
\paragraph{Official entry point}
We construct the directory layout expected by the authors' command-line pipeline, invoke their configuration generator and obfuscator, and attack the protected files they produce. The targets are 16 English conversations selected from a public LMSYS-Chat-1M mirror~\cite{lmsys}. The defender's magnitude thresholds are computed from the first eight target conversations; this setup is not an independent calibration/test split. Target plaintext is not supplied to the recovery routine. The saved per-file rates are 100\% for all 16 files. This record supports a controlled file-format test, not a production-service or remote-access claim. The historical harness reused an output directory and did not enforce subprocess success before reading files, so the saved summary alone does not certify a fresh end-to-end execution.

\paragraph{Multiple secret epochs}
Figure~\ref{fig:epochs} summarizes three models and five independently generated configurations per model. All heads calibrate in all 15 model--epoch cells. Mean recovery across epoch means is 100\% for Qwen2.5-0.5B, 99.64\% for Llama-3.2-1B, and 99.34\% for Phi-3-mini; the lowest model-level epoch mean is 99.29\%. Wrong-secret controls average 0.05\%, 0.01\%, and 0.16\%, respectively. The same ten public conversations recur across epochs and overlap earlier public-data experiments, so we treat this as secret-epoch stability, not held-out-session generalization. The result file retains rounded session rates rather than raw counts.

\begin{figure*}[t]
\centering
\includegraphics[width=.88\textwidth]{key_epochs.pdf}
\caption{Secret-epoch stability on three models. Every epoch calibrates all heads. The y-axis in (a) is deliberately zoomed; (b) uses demixers from the wrong configuration. The conversations are not an exclusive held-out set.}
\label{fig:epochs}
\end{figure*}

\paragraph{Lifecycle simulation}
In a controlled Qwen2.5-0.5B service simulation, one calibration request is followed by ten synthetic requests from different application domains. Recovery is 100\% on all ten. Replacing the complete secret configuration reduces a stale demixer to 0.78\%; one new calibration probe restores 100\% on the measured request. This result identifies the relevant unit of exposure: a shared secret epoch. It does not estimate tenant co-location prevalence in real deployments.

The lifecycle and epoch experiments support different interpretations of reuse. Recalibrating after rotation shows that a new shared epoch can again supply the same type of calibration opportunity; it does not mean the old demixer survives rotation. The separate request-refresh experiment asks whether an old demixer transfers when its row secrets are no longer shared. Reporting both observations makes the scope of the proposed intervention explicit.

\subsection{RQ4: bounded privacy consequence}
We evaluate closed-set identification on 30 synthetic cases. For each of six value types---SSN-like strings, card numbers, passwords, medical record numbers, names, and addresses---we generate a fixed list of 100 candidates and use the first five as truths. A target value is embedded in neutral context. The attacker aggregates recovered tokens from all complete blocks and all heads; it is not told which block contains the value. A candidate is scored by the fraction of its unique tokens present in the recovered multiset.

Let $g$ be the number of candidates with a strictly higher score than the truth and $m$ the size of the truth's tie group. The truth attains the maximum in all 30 cases ($g=0$); it is the unique maximum in 21/30. Under uniform random tie-breaking within the recorded group, expected top-1, top-5, and top-10 success are 77.36\%, 90.95\%, and 95.24\%, respectively. Uniform guessing from 100 candidates gives 1\%, 5\%, and 10\%. Figure~\ref{fig:privacy} shows the top-1 result by value type.

The expected top-$k$ contribution of a case is
\begin{equation}
 p_k(g,m)=\min\{1,\max\{0,(k-g)/m\}\}.
 \label{eq:ties}
\end{equation}
There are $g$ candidates ahead of the tie group and $m$ equally likely positions within it. If the remaining top-$k$ slots cover the entire group, success is one; otherwise the fraction of covered positions determines the expectation. Thus $g=0$ means that the truth is not outranked, whereas unique identification additionally requires $m=1$. We average these per-case expectations over the 30 retained cases, rather than treating tied maxima as deterministic successes.

The candidate score tests token membership, not order or a verified field span. A candidate can share tokens with the surrounding context or with another candidate. Even perfect recovery of a target multiset can therefore leave several values tied. Conversely, identifying a candidate from a supplied list does not imply that the method can generate that value without candidate knowledge. These distinctions explain why the privacy outcome and the main multiset metric are complementary rather than interchangeable.

\begin{figure}[t]
\centering
\includegraphics[width=\columnwidth]{privacy.pdf}
\caption{Closed-set identification in the 30-case synthetic study. Each bar contains five cases. Ties are broken uniformly in expectation.}
\label{fig:privacy}
\end{figure}

This experiment is evidence of consequence, not a population estimate. Truths are the first five generated candidates rather than uniform random indices; the displayed 1\% line is therefore a uniform-guess reference, not a matched optimal prior. Only the truth score and $(g,m)$ summaries were retained, not all candidate scores or recovered multisets. An ``entity complete'' flag used a token-set approximation rather than a verified character span. We consequently do not claim unconstrained credential recovery, stable per-type rates, or complete-field reconstruction.

\subsection{Shortening the secret lifetime}
The mechanism suggests refreshing the reused row transform $S$ and marker $a$ while leaving the fused column transform $M$ static. In one Qwen2.5-0.5B experiment, the saved same-epoch recovery rate is 97.44\%, compared with 0\% when a demixer calibrated for request A is applied after $S,a$ are regenerated for request B. The result JSON retains rounded rates rather than raw counts for this comparison. Generating the corresponding matrices for the full model takes 2.9\,ms in a CPU microbenchmark.

This is a mitigation candidate, not a validated repair. The saved round-trip check does not establish full K/V functional equivalence, and the test does not cover adaptive within-request calibration, shared prefixes, multi-turn identity, device-transfer cost, or end-to-end serving latency. The defensible conclusion is narrower: request-specific $S,a$ prevents the \emph{existing} cross-request demixer from transferring in the measured setting.

\subsection{Answers to the research questions}
The observed state counts and parameter errors support the predicted marker orbit (RQ1). Recovery is high across the tested models and configurations, while fresh-secret control recovery remains between 0 and 0.625\% (RQ2). The attack consumes the official protected-cache output and remains effective across generated secret epochs, provided calibration and target caches share the relevant secrets (RQ3). The synthetic closed-set study exceeds the uniform-guess reference, with candidate selection and incomplete traces limiting its interpretation. Shortening the row-secret lifetime interrupts the measured transfer, but functional correctness and general security remain unverified (RQ4).
